O potencial gravitacional Newtoniano a uma distância $r$ do centro da Terra é:
\[
V_g(r) = -\frac{G\,M_E\,m}{r}.
\]
À altitude $z$ acima da superfície, a distância ao centro é $r = R + z$ (onde $R$ é o raio terrestre). Substituindo:
\[
\boxed{V_g(z) = -\frac{G\,M_E\,m}{R+z}.}
\]
Para $z = 0$ (superfície), recupera-se o potencial de referência $V_g(0) = -GM_Em/R$. Para $z\to\infty$, $V_g \to 0$ (ponto de referência convencional no infinito).